% -----------------------------------------------------------
\section{Mystery}
% -----------------------------------------------------------
	\label{MYSTERY}
	
	See also \hyperref[GODLINESS]{Godliness},
	
% % ----------------------
% \subsection{See also}
% % ----------------------

% % ----------------------
% \subsection{1828 American Dictionary of the English Language} % *1828
% % ----------------------
	% \label{MYSTERY-1828}

	% \begin{compactitem}
		% \item
	% \end{compactitem}

% % ----------------------
% \subsection{Oxford English Dictionary} % *OED
% % ----------------------
	% \label{MYSTERY-OED}

	% \begin{compactitem}
		% \item[]
	% \end{compactitem}

% ----------------------
\subsection{Scriptural Usage} % *SCRIPTURES
% ----------------------
	\label{MYSTERY-SCRIPTURES}

	% % -----
	% \subsubsection{Old Testament} % *OT
	% % -----
		% \label{MYSTERY-OT}
	
		% % --
		% \paragraph{KJV}
		% % --
			% \label{MYSTERY-OT-KJV}
	
			% \begin{tabular}{ll}
				% Total occurrences in the OT & 
			% \end{tabular}
			
			% \begin{compactdesc}
				% \item[]
			% \end{compactdesc}
			
		% % --
		% \paragraph{Hebrew Bible}
		% % --
			% \label{MYSTERY-HEBREW}
		
			% %
			% \subparagraph{\texorpdfstring{\he{sample word}}{}} % paste actual hebrew text here
			% %
				% \label{PAGA} % whatever is in step bible without periods
			
				% \begin{compactdesc}
					% \item[HALOT , PDF ] \textbf{} % Use the 1994 PDF
						% \begin{compactitem}
							% \item[1]
						% \end{compactitem}
					% \item[BDB , PDF ] \textbf{}
						% \begin{compactitem}
							% \item[1]
						% \end{compactitem}
					% \item[DCH PDF ] \textbf{}
						% \begin{compactitem}
							% \item[1]
						% \end{compactitem}
				% \end{compactdesc}
				
				% \begin{compactdesc}
					% \item[Some OT uses] \textbf{}
						% \begin{compactdesc}
							% \item[]
						% \end{compactdesc}
				% \end{compactdesc}
			
		% % --
		% \paragraph{Septuagint}
		% % --
			% \label{MYSTERY-LXX}
		
			% %
			% \subparagraph{\texorpdfstring{\gr{sample word}}{}}
			% %
		
	% -----
	\subsubsection{New Testament} % *NT
	% -----
		\label{MYSTERY-NT}
	
		% --
		\paragraph{KJV}
		% --
			\label{MYSTERY-NT-KJV}
			
	
			\begin{tabular}{ll}
				Total occurrences in the NT & 27
			\end{tabular}
			
			\begin{compactdesc}
				\item[{\hyperref[MATTHEW-13-11]{Matthew 13:11}}] He answered and said unto them, Because it is given unto you to know the mysteries of the kingdom of heaven, but to them it is not given.
				\item[{\hyperref[MARK-4-11]{Mark 4:11}}] And he said unto them, Unto you it is given to know the mystery of the kingdom of God: but unto them that are without, all these things are done in parables:
				\item[{\hyperref[LUKE-8-10]{Luke 8:10}}] And he said, Unto you it is given to know the mysteries of the kingdom of God: but to others in parables; that seeing they might not see, and hearing they might not understand.
				\item[{\hyperref[ROMANS-11-25]{Romans 11:25}}] For I would not, brethren, that ye should be ignorant of this mystery, lest ye should be wise in your own conceits; that blindness in part is happened to Israel, until the fulness of the Gentiles be come in.
				\item[{\hyperref[ROMANS-16-25]{Romans 16:25}}] Now to him that is of power to stablish you according to my gospel, and the preaching of Jesus Christ, according to the revelation of the mystery, which was kept secret since the world began,
				\item[{\hyperref[1CORINTHIANS-2-7]{1 Corinthians 2:7}}] But we speak the wisdom of God in a mystery, even the hidden wisdom, which God ordained before the world unto our glory:
				\item[{\hyperref[1CORINTHIANS-4-1]{1 Corinthians 4:1}}] LET a man so account of us, as of the ministers of Christ, and stewards of the mysteries of God.
				\item[{\hyperref[1CORINTHIANS-13-2]{1 Corinthians 13:2}}] And though I have the gift of prophecy, and understand all mysteries, and all knowledge; and though I have all faith, so that I could remove mountains, and have not charity, I am nothing.
				\item[{\hyperref[1CORINTIHANS-14-2]{1 Corinthians 14:2}}] For he that speaketh in an unknown tongue speaketh not unto men, but unto God: for no man understandeth him; howbeit in the spirit he speaketh mysteries.
				\item[{\hyperref[1CORINTHIANS-15-51]{1 Corinthians 15:51}}] Behold, I shew you a mystery; We shall not all sleep, but we shall all be changed,
				\item[{\hyperref[EPHESIANS-1-9]{Ephesians 1:9}}] Having made known unto us the mystery of his will, according to his good pleasure which he hath purposed in himself:
				\item[{\hyperref[EPHESIANS-3-3]{Ephesians 3:3}}] How that by revelation he made known unto me the mystery; (as I wrote afore in few words,
				\item[{\hyperref[EPHESIANS-3-4]{Ephesians 3:4}}] Whereby, when ye read, ye may understand my knowledge in the mystery of Christ)
				\item[{\hyperref[EPHESIANS-3-9]{Ephesians 3:9}}] And to make all men see what is the fellowship of the mystery, which from the beginning of the world hath been hid in God, who created all things by Jesus Christ:
				\item[{\hyperref[EPHESIANS-5-32]{Ephesians 5:32}}] This is a great mystery: but I speak concerning Christ and the church.
				\item[{\hyperref[EPHESIANS-6-19]{Ephesians 6:19}}] And for me, that utterance may be given unto me, that I may open my mouth boldly, to make known the mystery of the gospel,
				\item[{\hyperref[COLOSSIANS-1-26]{Colossians 1:26}}] Even the mystery which hath been hid from ages and from generations, but now is made manifest to his saints:
				\item[{\hyperref[COLOSSIANS-1-27]{Colossians 1:27}}] To whom God would make known what is the riches of the glory of this mystery among the Gentiles; which is Christ in you, the hope of glory:
				\item[{\hyperref[COLOSSIANS-2-2]{Colossians 2:2}}] That their hearts might be comforted, being knit together in love, and unto all riches of the full assurance of understanding, to the acknowledgement of the mystery of God, and of the Father, and of Christ;
				\item[{\hyperref[COLOSSIANS-4-3]{Colossians 4:3}}] Withal praying also for us, that God would open unto us a door of utterance, to speak the mystery of Christ, for which I am also in bonds:
				\item[{\hyperref[2THESSALONIANS-2-7]{2 Thessalonians 2:7}}] For the mystery of iniquity doth already work: only he who now letteth will let, until he be taken out of the way.
				\item[{\hyperref[1TIMOTHY-3-9]{1 Timothy 3:9}}] Holding the mystery of the faith in a pure conscience.
				\item[{\hyperref[1TIMOTHY-3-16]{1 Timothy 3:16}}] And without controversy great is the mystery of godliness: God was manifest in the flesh, justified in the Spirit, seen of angels, preached unto the Gentiles, believed on in the world, received up into glory.
				\item[{\hyperref[REVELATION-1-20]{Revelation 1:20}}] The mystery of the seven stars which thou sawest in my right hand, and the seven golden candlesticks.  The seven stars are the angels of the seven churches: and the seven candlesticks which thou sawest are the seven churches.
				\item[{\hyperref[REVELATION-10-7]{Revelation 10:7}}] But in the days of the voice of the seventh angel, when he shall begin to sound, the mystery of God should be finished, as he hath declared to his servants the prophets.
				\item[{\hyperref[REVELATION-17-5]{Revelation 17:5}}] And upon her forehead was a name written, MYSTERY, BABYLON THE GREAT, THE MOTHER OF HARLOTS AND ABOMINATIONS OF THE EARTH.
				\item[{\hyperref[REVELATION-17-7]{Revelation 17:7}}] And the angel said unto me, Wherefore didst thou marvel?  I will tell thee the mystery of the woman, and of the beast that carrieth her, which hath the seven heads and ten horns.
			\end{compactdesc}
			
		% --
		\paragraph{Greek New Testament} % *GREEK
		% --
			\label{MYSTERY-GREEK}
		
			%
			\subparagraph{\texorpdfstring{\gr{μυστήριον}}{}}
			%
				\label{MUSTERION}
			
				\begin{compactdesc}
					\item[BDAG 586b] `secret, secret rite, secret teaching, mystery' a religious technical term (predominantly plural) applied in the Greco-Roman world mostly to the mysteries with their secret teachings, religious and political in nature, concealed within many strange customs and ceremonies. The principal rites remain unknown because of a reluctance in antiquity to divulge them.
						\begin{compactitem}
							\item[1] \textbf{the unmanifested or private counsel of God, \textit{(God's)} secret,} the secret thoughts, plans, and dispensations of God
							\item[1a] In the gospels \gr{μυστήριον} is found only in one context, where Jesus says to the disciples who have asked for an explanation of the parable(s)
							\item[1b] The Pauline literature has \gr{μυστήριον} in 21 places. A \textit{secret} or \textit{mystery}, too profound for human ingenuity, is God's reason for the partial hardening of Israel's heart (\hyperref[ROMANS-11-25]{Romans 11:25})...
							\item[1c] In Revelation \gr{μυστήριον} is used in reference to the mysterious things portrayed there.
							\item[2] \textbf{that which transcends normal understanding, \textit{transcendent/ultimate reality, secret}}, with focus on Israelite/Christian experience
							\item[2a] \hyperref[1TIMOTHY]{1 Timothy} uses \gr{μυστήριον} as a formula...\gr{τὸ τῆς εὐσεβείας μυστήριον}, \textit{the secret of (our) piety} (\hyperref[1TIMOTHY-3-16]{1 Timothy 3:16}) \textit{(See below for this verse in Greek.)}
							\item[2b] In Ignatius: the death and resurrection of Jesus as \gr{μυστήριον}
							\item[2c] Quite difficult is the saying about the tried and true prophet...
							\item[2d] \gr{μυστήριον} offucs often in Diognetus...
						\end{compactitem}
					\item[CGL 953a, PDF 993] \textbf{}
						\begin{compactitem}
							\item[1] mystery or secret rites, \textbf{mysteries} (Heraclitus, Herodotus, Plutarch)
							\item[2] \textbf{Mysteries} (especially of Demeter and Kore at Eleusis)
							\item[3] (figurative, referring to arcane teachings or philosophical doctrines) \textbf{mysteries}
							\item[4] \textbf{mysteries, secrets} (with genitive of the Kingdom of Heaven) New Testament (also singular)
						\end{compactitem}
					\item[LSJ 1156b, PDF 1227] \textbf{}
						\begin{compactitem}
							\item[I.1] \textit{mystery} or \textit{secret rite}
							\item[I.2a] \textit{mystic implements} and \textit{ornaments}
							\item[I.2b] later, \textit{object used in magical rites, talisman}
							\item[I.3] generally, \textit{mystery, secret}
							\item[I.4] \textit{secret revealed by God}, i.e. \textit{religious} or \textit{mystical truth} \textit{(This is the sense applicable to the New Testament.)}
							\item[I.5] \textit{superstition}
							\item[II] Dionysius the tyrant called \textit{mouse-holes} \gr{μυστήρια}
						\end{compactitem}
				\end{compactdesc}
				
				\begin{compactdesc}
					\item[Some NT uses] \textbf{}
						\begin{compactdesc}
							\item[Matthew 13:11]
							\item[Mark 4:11] 
							\item[Luke 8:10]
							\item[Romans 11:25]
							\item[Romans 16:25] 
							\item[1 Corinthians 4:1]
							\item[1 Corinthians 15:51]
							\item[Ephesians 3:3]
							\item[Ephesians 3:4]
							\item[Ephesians 3:9]
							\item[2 Thessalonians 2:7] 
							\item[1 Timothy 3:16] \gr{καὶ ὁμολογουμένως μέγα ἐστὶν τὸ τῆς εὐσεβείας μυστήριον·  Ὃς ἐφανερώθη ἐν σαρκί, ἐδικαιώθη ἐν πνεύματι, ὤφθη ἀγγέλοις, ἐκηρύχθη ἐν ἔθνεσιν, ἐπιστεύθη ἐν κόσμῳ, ἀνελήμφθη ἐν δόξῃ.}
							\item[Revelation 17:7]
						\end{compactdesc}
				\end{compactdesc}
		
	% % -----
	% \subsubsection{Book of Mormon} % *BOM
	% % -----
		% \label{MYSTERY-BOM}
	
		% \begin{tabular}{ll}
			% Total occurrences in the BOM & 
		% \end{tabular}
		
		% \begin{compactdesc}
			% \item[]
		% \end{compactdesc}
		
	% -----
	\subsubsection{Doctrine and Covenants} % *DC
	% -----
		\label{MYSTERY-DC}
	
		\begin{tabular}{ll}
			Total occurrences in the DC & 25
		\end{tabular}
		
		\begin{compactdesc}
			\item[{\hyperref[DC19-8]{DC 19:8}}]
			\item[{\hyperref[DC19-10]{DC 19:10}}]
		\end{compactdesc}
		
	% % -----
	% \subsubsection{Pearl of Great Price} % *PGP
	% % -----
		% \label{MYSTERY-PGP}
	
		% \begin{tabular}{ll}
			% Total occurrences in the PGP & 
		% \end{tabular}
		
		% \begin{compactdesc}
			% \item[]
		% \end{compactdesc}
		
% ----------------------
\subsection{Key Phrases} % *KEYPHRASES *PHRASES
% ----------------------
	\label{MYSTERY-PHRASES}

	\begin{compactdesc}
		\item[mystery of godliness] (See this phrase and commentary \hyperref[GODLINESS-PHRASES]{here}, at \hyperref[GODLINESS]{Godliness}.)
			% \begin{compactdesc}
				% \item[Bible anchor verses] \phantom{.}
					% \begin{compactdesc}
						% \item[Romans 5:5] \gr{}
					% \end{compactdesc}
				% \item[Commentary] ---
			% \end{compactdesc}
	\end{compactdesc}
	
% % ----------------------
% \subsection{Bible Dictionaries} % *DICTIONARIES
% % ----------------------
	% \label{MYSTERY-BD}

% % ----------------------
% \subsection{Restoration Usage} % *RESTORATION
% % ----------------------
	% \label{MYSTERY-RESTORATION}

	% % -----
	% \subsubsection{Joseph Smith} % *JS
	% % -----
		% \label{MYSTERY-JS}
	
		% \begin{compactdesc}
			% \item[{\hyperref[KEY]{Date/Source}}]
		% \end{compactdesc}
		
	% % -----
	% \subsubsection{Temple Ordinances}
	% % -----
		% \label{MYSTERY-TEMPLE}
	
		% \begin{compactdesc}
			% \item[{\hyperref[KEY]{Date/Source}}]
		% \end{compactdesc}
	
	% % -----
	% \subsubsection{Brigham Young} % *BY
	% % -----
		% \label{MYSTERY-BY}
	
		% \begin{compactdesc}
			% \item[{\hyperref[KEY]{Date/Source}}]
		% \end{compactdesc}
	
% % ----------------------
% \subsection{Commentary and Studies} % *COMMENTARY *STUDIES
% % ----------------------
	% \label{MYSTERY-COMMENTARY}

	% % -----
	% \subsubsection{Key Points}
	% % -----
		% \label{MYSTERY-KEYS}
	
		% \begin{compactitem}
			% \item
		% \end{compactitem}
		
	% % -----
	% \subsubsection{Sample Study}
	% % -----
		% \label{MYSTERY-study}